\chapter*{Introduction} % chapter* je necislovana kapitola
\addcontentsline{toc}{chapter}{Introduction} % rucne pridanie do obsahu
\markboth{Introduction}{Introduction} % vyriesenie hlaviciek

Modern neural networks contain hundreds of millions to hundreds of billions of parameters. Most of the computational cost of running them at inference is spent in linear layers, multiplying activations by large dense weight matrices. A natural way to reduce this cost is \emph{pruning}: setting a fraction of the weights to zero, ideally without losing model quality. If the pruning pattern matches what hardware can skip, the savings translate directly into faster inference.

A simple type of pruning is \emph{unstructured pruning}, which zeros out individual weights based on their importance to the model; the resulting sparse matrix is irregular and hard for hardware to accelerate. A potential solution to this problem is \emph{block pruning}, which zeros out fixed-shape non-overlapping blocks of weights, producing a sparsity pattern that maps directly onto specialized matrix multiplication routines on modern GPUs. The trade-off is that block pruning is more destructive: by forcing groups of weights to be removed together, it sometimes removes important weights that an unstructured pruner would have kept. The wider the block, the more enhanced this effect.

The TETRIS algorithm \cite{tetris} reduces this loss by applying a permutation along each dimension of the weight tensor before pruning. Reordering elements along a dimension does not change the network's output\footnote{In our implementation the permutation is used only as a mask-finding device: we permute the weights, apply the mask, and apply the inverse permutation before writing the weights back, so the model itself is never reordered. See the beginning of Chapter~\ref{chap:methods} for details.}, but it changes which weights end up grouped into the same block, and therefore which weights get pruned. The Original TETRIS searches the space of permutations greedily, optimizing one dimension at a time and swapping pairs of indices within that dimension whenever the swap reduces the pruned mass, stopping at the first local optimum.

In this thesis we propose \emph{Global Assignment TETRIS} (GA-TETRIS), a variant of TETRIS that replaces the greedy single-swap search with an exact solution to the underlying assignment problem. We observe that for a fixed pruning mask, finding the optimal permutation reduces to minimum-weight perfect matching in a bipartite graph, which can be solved in polynomial time using the linear assignment algorithm. Iterating this exact step --- alongside noise injection and a second phase of random-swap refinement to escape local optima --- forms the GA-TETRIS algorithm.

We compare GA-TETRIS against the Original TETRIS, three simpler heuristics we introduce (Sort-by-Norm, Random-Swaps, and Random-Swaps with sorted initialization), and the no-permutation Block-Wanda baseline. We evaluate on three pretrained models: a Small ViT (0.9M parameters), a Larger ViT (13.4M), and SmolLM2-360M. The comparison is conducted at two levels: a per-layer view of how well each algorithm optimizes the pruning objective directly, and a whole-model view of what each algorithm does to downstream task quality.

The findings are mixed. On the per-layer objective, GA-TETRIS produces the largest score improvements at wide block widths consistently across all three models, with the algorithm differences concentrated in early transformer blocks. On whole-model task metrics, however, all algorithms collapse rapidly as block width grows: only the Larger ViT at block $1 \times 2$ retains substantial accuracy. In that single regime, GA-TETRIS is the strongest algorithm by 1.85 percentage points. On SmolLM2 at block $1 \times 2$, simple Sort-by-Norm produces the lowest perplexity, though still roughly seven times worse than the unpruned baseline. The best algorithm depends on the model.
